\documentclass[11pt]{article}
\usepackage[margin=1in]{geometry}
\usepackage{amsmath,amssymb,amsthm,booktabs,microtype,hyperref}
\hypersetup{hidelinks,pdfauthor={T},pdftitle={Robust Consultation for Option-Preserving Assistance}}
\newtheorem{proposition}{Proposition}
\title{Robust Consultation for Option-Preserving Assistance}
\author{T}
\date{}
\begin{document}
\maketitle

\begin{abstract}
Acting immediately can consume an option that a beneficiary would
have preferred to keep. In a two-stage decision model, preserving
a menu and learning the beneficiary's state yields a value
of information with an explicit cost threshold. When both the
beneficiary's likely preference and the reliability of a reply
are uncertain, we derive the maximin rule for rectangular uncertainty.
When the two uncertainties are coupled, mixing immediate assistance
with consultation can strictly improve the guarantee. A finite
exact-arithmetic audit checks both settings. For any finite
joint scenario set, a two-dimensional distance formula gives
the exact optimal guarantee. For a finite batch, we distinguish a
hard question quota from its average-rate counterpart. When the
scenario can respond to the remaining quota, a backward recursion
prices the loss of one question slot.
\end{abstract}

\section*{Summary and scope}
\paragraph{Question.} When is it worth asking a beneficiary
which already permitted option to use?
\paragraph{Finding.} The value depends on the joint uncertainty
about preferences and reply quality. Its convex hull gives a
sharp guarantee, and randomizing between asking and acting can
improve that guarantee.
\paragraph{Evidence and scope.} Exact checks cover
rectangular and coupled models. The fitted comparison uses
synthetic costs and labels. A reply can inform the choice
within the authorized menu, and permission is checked separately.

\section{Why keeping an option changes the decision}
An assistant can offer two possible ways to resolve a problem,
prepare either one, or commit to one immediately. Suppose permitted
actions form the same finite menu $\mathcal A$ in each state
$\theta$. Acting immediately yields
\[
 V_0=\max_{a\in\mathcal A}E[u(a,\theta)].
\]
If a reversible step preserves the menu and a query perfectly
reveals $\theta$ before action, the expected value after paying
query cost $c$ is $V_1=E[\max_a u(a,\theta)]-c$.
This is a value-of-information comparison in the tradition of
Blackwell's order \cite{rauh}, with an explicit cost for the more
informative experiment. An irreversible operation can remove later
options, so the calculation assumes a preserved menu.

\begin{proposition}
The reversible and informed route strictly improves expected
utility exactly when
$c<E[\max_a u(a,\theta)]-\max_a E[u(a,\theta)]$.
The right side is nonnegative.
\end{proposition}
\begin{proof}
Subtract $V_0$ from $V_1$. For every fixed action the statewise
maximum is at least its utility, so expectation dominates the
best fixed action. Equality holds whenever the same action
maximizes in every positive-probability state.
\end{proof}

\section{An exact threshold}
Take two equally likely states. A tailored choice yields $4$
and the wrong choice yields $-2$. An immediate commitment gives
$1$, while perfect information yields $4$ before cost. The
threshold is $3$. At cost $2$ the reversible route is better,
and at cost $4$ acting immediately has higher expected utility
within this toy contract. With a symmetric answer channel of
accuracy $\rho\geq1/2$, using the indicated option yields
$6\rho-2$. Its improvement before cost is $6\rho-3$. A
bare instruction to always ask would lose in several of these
cases.

\section{When consultation reliability is uncertain}
Suppose both actions are permitted and exactly one matches the
beneficiary's state. Normalize the matching and mismatching payoffs
to one and zero, respectively, and write $d\ge0$ for consultation
cost in those units. Let $p$ be the probability that option one
matches. The binary answer is correct with the same probability
$r$ in either state. A scenario $(p,r)$ is fixed before the policy
draws any private random coin. The policy knows the uncertainty
set, and the reply is its only observation after choosing to ask.

The four relevant deterministic payoffs are
\[
 p,\qquad 1-p,\qquad r-d,\qquad 1-r-d.
\]
They correspond to choosing option one immediately, choosing
option zero immediately, following the reply, and reversing the
reply. A policy that pays to ask and then always chooses the same
option is weakly dominated by its immediate counterpart.

\begin{proposition}
For an uncertainty rectangle $[\ell,u]\times[r_-,r_+]$, the
optimal worst-case expected payoff is
\[
 \max\bigl\{\ell,1-u,\tfrac12,r_--d,1-r_+-d\bigr\}.
\]
\end{proposition}
\begin{proof}
The first two values are attained by the immediate options. An
equal mixture of those options attains $1/2$ in every scenario.
Following or reversing the reply attains the last two values,
respectively. For an upper bound choose the admissible scenario
whose $p$ and $r$ are each closest to $1/2$. At that scenario
the best immediate payoff is $\max\{\ell,1-u,1/2\}$ and the
best query payoff is $\max\{r_--d,1-r_+-d,1/2-d\}$. Averaging
among policies cannot exceed the highest payoff there. Since
$d\ge0$, the claimed formula follows.
\end{proof}

The rectangle can be too pessimistic if the weak regimes for
immediate action and consultation cannot occur together. For
example, at $d=1/10$ take the two scenarios
$(p,r)=(19/20,11/20)$ and $(11/20,19/20)$. The best pure
policy guarantees $11/20$. An equal mixture of immediate
option one and following the reply gives $7/10$ in both scenarios.
To certify optimality, give the two scenarios weights $3/8$ and
$5/8$. Each of those two policies then has mean $7/10$, and
the other two have smaller means. Every policy's worst-case
payoff is bounded above by its mean under these weights.
The enclosing rectangle gives only $11/20$, as it permits
low values of both $p$ and $r$ at once.

\section{The exact value of a joint scenario set}
Let $S$ be any finite nonempty set of pairs $(p,r)$ under the
same symmetric answer and cost assumptions. For a fixed
policy, its expected payoff is affine in these two coordinates.
Consequently the worst payoff over $S$ agrees with the worst
payoff over its convex hull.

\begin{proposition}
The optimal worst-case expected payoff is
\[
\frac12+\min_{(p,r)\in\operatorname{conv}S}
  \max\left\{\left|p-\frac12\right|,
              \left|r-\frac12\right|-d,\ 0\right\}.
\]
\end{proposition}
\begin{proof}
The finite game has four pure policies and the scenarios in
$S$ as its opposing choices. A randomized policy is a probability
vector over the four policies. Linear-programming duality
equates its best worst-case payoff to the smallest best-policy
payoff under a mixture of scenarios. The mixture is represented
by its mean pair $(p,r)\in\operatorname{conv}S$ because every
pure payoff is affine. The maximum of the four pure payoffs
at that pair is
$\max\{p,1-p,r-d,1-r-d\}$. Subtracting $1/2$ gives the
displayed distance expression.
\end{proof}

Geometrically, expand the rectangle
$|p-1/2|\le\gamma$, $|r-1/2|\le d+\gamma$
from $\gamma=0$ until it meets the hull. The first intersection
gives the guaranteed advantage $\gamma$ over a fair immediate
choice. Inspecting individual scenarios can miss the least
favorable mixture in the interior. For example, the triangle
with vertices $(3/4,1/2),(1/4,3/4),(1/4,1/4)$ contains
$(1/2,1/2)$ in its interior. At zero query cost the robust
value is $1/2$, although checking only its three edges gives
a larger value. On 96 three- and four-scenario fixtures the
hull calculation agreed with independent enumeration of the
policy linear program's vertices. Mixed policies strictly
improved on pure policies in 78 fixtures. The counts concern
the declared synthetic grid.

In a consent-based workflow, the alternatives could be an
immediately usable, reversible draft and a draft tailored after a
short question. Joint measurements of preference prevalence and
reply reliability can support a scenario set. Use the rectangle
when the evidence supports only marginal ranges. A stronger
guarantee for a smaller joint set depends on that set covering
plausible cases. These policies neither authorize new actions nor
guarantee the welfare of each beneficiary.

\section{A limit on how often to ask}
A service may want a guarantee without interrupting every person.
Let $q\in[0,1]$ bound the probability of asking. This is distinct
from $d$, the utility cost each time the service asks. On the
scenario hull write $A(p)=\max(p,1-p)$ for the best immediate
value and $B(r)=\max(r,1-r)-d$ for the best query value.

\begin{proposition}
With a question rate at most $q$, the optimal worst-case payoff is
\[
 V(q)=\min_{(p,r)\in\operatorname{conv}S}
 \bigl[A(p)+q\max(0,B(r)-A(p))\bigr].
\]
A target $v>V(1)$ is infeasible. For $v\le V(1)$ the smallest
permitted rate is
\[
 \sup_{(p,r)\in\operatorname{conv}S:\ A(p)<v}
       \frac{v-A(p)}{B(r)-A(p)},
\]
with value zero if the set is empty.
\end{proposition}
\begin{proof}
At a fixed scenario mixture, let $x\le q$ be total probability
assigned to query policies. The best payoff at that $x$ is
$(1-x)A(p)+xB(r)$, since the two kinds of actions can be chosen
separately. Its maximum over $x$ is the stated expression.
The capped policy set is a compact polytope. The same finite-game
duality as before interchanges maximization over its vertices
and minimization over the scenario hull. If $A(p)\ge v$, the
point imposes no lower bound on $q$. Otherwise, feasibility at
$q=1$ implies $B(r)\ge v>A(p)$, and the expression is at least
$v$ exactly when $q\ge(v-A(p))/(B(r)-A(p))$. Every point must
satisfy its inequality, which yields the supremum.
\end{proof}

The expression in the minimization is the maximum of six affine
functions. Two represent immediate choices. Four represent asking
with probability $q$ and choosing one immediate action and one
reply rule. Thus the computation can use the vertices where
active affine functions change across the hull. A separate policy
oracle instead enumerates mixtures of the original four choices
subject to the explicit rate constraint.

In the earlier pair of scenarios with query cost $1/10$, a
one-quarter rate gives $5/8$ instead of the unrestricted $7/10$.
To guarantee $13/20$ requires a rate of one third. A rate above
one half gives no additional gain in this case. Across 256 exact
comparisons, the capped hull result agreed with the policy oracle.
A one-quarter cap strictly reduced the unrestricted value on 14
of 64 scenario sets.


\section{A hard quota for a finite batch}
Suppose $n$ cases share a quota of $k\le n$ questions. Each case has
the same authorized menu, cost and scenario set $S$. The objective
is the sum of case utilities. A scenario is fixed before the policy
draws its current coins. An oblivious scenario sequence is fixed
before the batch. A predictable adversary may instead choose each
scenario after observing the remaining quota. It has no access to
the policy's future coins. Neither party has a persistent hidden
state that constrains the future scenario set.

\begin{proposition}
For oblivious scenarios the optimal per-case guarantee under the
hard quota is $V(k/n)$. For predictable scenarios let $G(n,k)$
be the optimal total guarantee, with $G(0,k)=0$,
$G(n,0)=nV(0)$, and $G(n,k)=G(n,n)$ for $k>n$. Define
\[
 U(c)=\min_{(p,r)\in\operatorname{conv}S}
              \max\{A(p),\max(r,1-r)-c\}.
\]
For $n,k\ge1$,
\[
 G(n,k)=G(n-1,k)+U\bigl(d+G(n-1,k)-G(n-1,k-1)\bigr).
\]
\end{proposition}
\begin{proof}
For a fixed sequence, take a rate-capped optimal policy with question
probability $x\le k/n$. Pick $k$ distinct slots uniformly. For $k>0$,
ask independently in each eligible slot with probability $nx/k$.
The conditional action mixtures match the optimal one-case policy
and each slot has the intended marginal. This delivers $V(k/n)$
while respecting the quota on every path. With $k=0$, always use
the immediate mixture. For the converse, draw the scenarios
independently from a least favorable mixture attaining $V(k/n)$.
Past observations cannot predict a fresh scenario draw.
If the expected fraction of questions is $y\le k/n$, the per-case
payoff under the mixture is at most $(1-y)A+yB$. Its maximum over
$y\le k/n$ equals $V(k/n)$. Some fixed sequence therefore gives
at most this value.

In the predictable game, retaining a slot leaves continuation
$G(n-1,k)$. Spending it leaves $G(n-1,k-1)$. The difference is
nonnegative by inclusion of feasible policies. Conditional on the
public history, all future scenarios remain possible. Subtracting
the common continuation therefore raises the effective question
cost by that difference. The one-case minimax identity gives $U$
at the raised cost. Backward induction supplies strategies for
both players at every quota.
\end{proof}

For the earlier two-scenario example, put $Z_n\sim\mathrm{Bin}(n,1/2)$.
The predictable value has the exact form
\[
 G(n,k)=\frac{11n}{20}+\frac3{10} E[\min(k,Z_n)].
\]
To see this, the expression's marginal value of a slot is
$(3/10)P(Z_n\ge k)$. For any such continuation charge $\Delta$
between zero and $3/10$, immediate option one and following the
reply dominate their counterparts. Their worst-case mixture is
optimized by asking with probability $1/2$ when a slot remains.
The two scenario payoffs meet there and the stage value is
$7/10-\Delta/2$. The recurrence for $E[\min(k,Z_n)]$ under one
additional fair coin yields the displayed value and the stated
boundary cases.

With two cases and one question, the oblivious guarantee is $7/10$
per case. The predictable guarantee is $53/80$. At $k=n/2$ for
even $n$, the total gap is $(3/20)E|Z_n-n/2|$, so the gap per case
is at most $3/(40\sqrt n)$. This follows from symmetry and the
binomial variance. The verifier compared 336 reachable states
using two exact stage solvers, checked the binomial expression at
560 states, and checked the two-case saddle against all eight
scenario trees and twelve pure plans.

A support team with a fixed daily question allowance can randomly
reserve slots if its case mix is fixed before those choices.
If requests arrive in response to disclosed capacity, use a
state-dependent guarantee. Neither calculation changes the
permission needed to act within an individual case.

\section{Building examples that transfer}
Generate episodes with identical first messages and different
beneficiary preferences. Include a reversible preparatory step,
an immediate commitment, and a short question with an explicit
delay or attention cost. In independent test families measure
whether the model retains meaningful control for the affected
person and whether it delivers helpful work when a question has
little value. Permission is a separate constraint on the menu,
including any cost of delay in time-sensitive cases. Collect
preferences under consent and keep emergency scenarios subject
to appropriate authority.

\newpage
\section{Verification and limits}
The checker compares the statewise formula against explicit
enumeration of four contingent policies across all two-by-two
tables over $\{-2,0,4\}$ and five priors. All 405 agree.
An exact second audit checked the robust value on 270 rectangles
and compared primal policy mixtures with dual scenario bounds on
315 scenario-pair and cost combinations. Mixing strictly improved
on the best pure policy in 118 cases. These counts describe the fixed rational
grid. The proposal depends on realistic costs and faithful
preferences. A falsely reversible operation, an answer channel
with state-dependent error, or an omitted stakeholder can
invalidate the model.
\section{Training and testing when to ask}
The learned logistic policy trained on 320 synthetic examples per
seed using the rule \(6\rho-3>c\). It chose whether to pay for
information. Across 3,200 fresh episodes, it queried 1,690 times.
It achieved mean realized utility 1.512, compared with 1.094 for
always asking and 0.976 for never asking. All policies shared
the same realized consultation and immediate-action outcomes
for each episode. The per-seed improvement over never asking
ranged from 0.199 to 0.857 utility units. The evaluation depends
on an available, reversible menu and the simulated accuracy
and cost distribution. It supplies an empirical check on the
decision rule in that environment.

\paragraph{Review status.} Independent manuscript review and a language-model study remain open.

\begin{thebibliography}{9}
\bibitem{rauh} Johannes Rauh, Pradeep Kr.~Banerjee, Eckehard
Olbrich, J\"urgen Jost, Nils Bertschinger, and David H.~Wolpert.
Coarse-Graining and the Blackwell Order. \emph{Entropy}
19(10), 527, 2017. \url{https://arxiv.org/pdf/1701.07602}.
\end{thebibliography}
\end{document}
